\section{Results}
\label{sec:results}

We evaluate on three workloads: \textsc{tpch-q3}, a synthetic uniform
join, and a skewed join with a Zipfian key distribution. All runs use a
single core with a 32~MiB last-level cache, and all numbers are the median
of eleven runs.

\begin{table}[htbp]
\centering
\begin{tabular}{lrrr}
\toprule
Workload & Arrival order & Grouped & Gain \\
\midrule
\textsc{tpch-q3} & 1.00 & 1.71 & \(1.71\times\) \\
Uniform & 1.00 & 1.66 & \(1.66\times\) \\
Zipfian & 1.00 & 1.79 & \(1.79\times\) \\
\bottomrule
\end{tabular}
\caption{Normalized throughput by workload, buffer \(n = 32\).}
\label{tab:main}
\end{table}

The gain is largest on the skewed workload, which is the case the
residency model predicts: a skewed key distribution concentrates probes
on few partitions, so grouping keeps those partitions resident across
many more batches than a uniform mix does.

\subsection{Where the model is wrong}

The model assumes a partition is either resident or absent. Real caches
hold partial partitions, so the predicted miss rate is pessimistic by
roughly four points across every workload we measured; the predicted
\emph{ordering} of configurations was correct in every case, which is
what the scheduler actually consumes~\cite{blanas2011}.
